\chapter{Movement and Items}

본 장에서는 실제 기보를 빠르게 계산하기 위한 최소한의 상태 언어와 아이템의 효과를 정리한다. 목적은 주사위의 다음 윗면, 이동거리, 도달 가능 칸을 실전에서 정확하게 읽는 것이다.

\section{Gesture States and Roll Analysis}

RPSC의 주사위는 실제로 24가지 exact orientation을 가진다. 정확한 기보 복원에서는 이 24가지 방향을 구별하지만, 윗면 변화와 이동거리를 빠르게 계산할 때에는 S/R/P가 세 공간축 가운데 어디에 놓였는지만 알면 충분하다.

\medskip
\noindent\textbf{Gesture State.}

가위, 바위, 보를 각각 $\mathrm{S}$, $\mathrm{R}$, $\mathrm{P}$로 나타낸다. 보드에 수직인 축을 $\mathrm{UD}$, North--South 축을 $\mathrm{NS}$, East--West 축을 $\mathrm{EW}$라 한다. Gesture State를
\[
    s=\langle a,b,c\rangle,
    \qquad
    \text{coordinate order }\langle \mathrm{UD},\mathrm{NS},\mathrm{EW}\rangle
\]
로 적는다. 첫 번째 성분 $a$가 현재 Top Gesture이다.

실제 판에서는 Top Gesture와 Wrist Axis만 보면 Gesture State를 바로 읽을 수 있다.

\begin{center}
    \renewcommand{\arraystretch}{1.16}
    \begin{tabular}{@{}ccc@{}}
        \toprule
        Top Gesture & Wrist Axis $\mathrm{NS}$ & Wrist Axis $\mathrm{EW}$ \\
        \midrule
        $\mathrm{S}$ & $\langle \mathrm{S},\mathrm{R},\mathrm{P}\rangle$ & $\langle \mathrm{S},\mathrm{P},\mathrm{R}\rangle$ \\
        $\mathrm{R}$ & $\langle \mathrm{R},\mathrm{S},\mathrm{P}\rangle$ & $\langle \mathrm{R},\mathrm{P},\mathrm{S}\rangle$ \\
        $\mathrm{P}$ & $\langle \mathrm{P},\mathrm{R},\mathrm{S}\rangle$ & $\langle \mathrm{P},\mathrm{S},\mathrm{R}\rangle$ \\
        \bottomrule
    \end{tabular}
\end{center}

Gesture State $s=\langle a,b,c\rangle$의 Base Roll Length를 $\ell(s)$로 적는다.
\[
    \ell(\langle a,b,c\rangle)\coloneqq
    \begin{cases}
        3,&a=\mathrm{S},\\
        4,&a=\mathrm{R},\\
        5,&a=\mathrm{P}.
    \end{cases}
\]

\medskip
\noindent\textbf{State Changes.}

세 공간축의 quarter-turn이 Gesture State에 만드는 변화만 다음 세 연산으로 적는다.
\[
\begin{aligned}
    \tau_{\mathrm{EW}}(\langle a,b,c\rangle)&\coloneqq\langle b,a,c\rangle,\\
    \tau_{\mathrm{NS}}(\langle a,b,c\rangle)&\coloneqq\langle c,b,a\rangle,\\
    \tau_{\mathrm{UD}}(\langle a,b,c\rangle)&\coloneqq\langle a,c,b\rangle.
\end{aligned}
\]

실제 Roll과 Rotation에는 다음처럼 적용한다.

\begin{center}
    \renewcommand{\arraystretch}{1.16}
    \begin{tabular}{@{}cc@{}}
        \toprule
        행동 & Gesture-State 변화 \\
        \midrule
        North / South Roll & $\tau_{\mathrm{EW}}$ \\
        East / West Roll & $\tau_{\mathrm{NS}}$ \\
        \texttt{RoN} / \texttt{RoS} & $\tau_{\mathrm{EW}}$ \\
        \texttt{RoE} / \texttt{RoW} & $\tau_{\mathrm{NS}}$ \\
        \texttt{RoL} / \texttt{RoR} & $\tau_{\mathrm{UD}}$ \\
        \bottomrule
    \end{tabular}
\end{center}

서로 반대인 두 Rotation은 Gesture State에서는 같은 변화를 만들지만 exact orientation에서는 서로 다르다. 따라서 실제 기보에서는 여섯 Rotation을 모두 구별한다.

\medskip
\noindent\textbf{Roll Word.}

한 이동의 Roll 방향을 시간 순서대로 적은 것을 Roll Word라 한다. 예를 들어
\[
    \texttt{W2: c1-c2-d2-d3-e3}
    \qquad\Longrightarrow\qquad
    NENE
\]
이다. 왼쪽부터 차례대로 $N,S$에는 $\tau_{\mathrm{EW}}$, $E,W$에는 $\tau_{\mathrm{NS}}$를 적용하면 최종 Top Gesture와 다음 Gesture State를 얻는다. Roll의 순서는 중요하므로 방향별 횟수만 세어서는 충분하지 않다.

\medskip
\noindent\textbf{Items.}

\begin{center}
    \renewcommand{\arraystretch}{1.18}
    \begin{tabular}{@{}L{0.17\textwidth}L{0.71\textwidth}@{}}
        \toprule
        아이템 & 실전 계산 \\
        \midrule
        Push & Gesture State는 바뀌지 않는다. 한 칸 민 뒤 $\ell(s)$번 Roll하며, 첫 Roll에서 Push 이전 칸으로 되돌아갈 수 없다. \\
        Rotation & 먼저 선택한 Rotation의 $\tau$를 적용하여 $s'$를 구한 뒤, \textbf{$\ell(s')$번} Roll한다. Rotation으로 Top Gesture와 이동거리가 바뀔 수 있다. \\
        Step Short & Gesture State는 그대로 두고 Roll 수를 $\ell(s)-1$로 정한다. \\
        Step Long & Gesture State는 그대로 두고 Roll 수를 $\ell(s)+1$로 정한다. \\
        \bottomrule
    \end{tabular}
\end{center}

예를 들어 초기 W2는
\[
    s=\langle \mathrm{R},\mathrm{S},\mathrm{P}\rangle
\]
이다. 이때 \texttt{RoN} 또는 \texttt{RoS}를 사용하면
\[
    \tau_{\mathrm{EW}}(s)=\langle \mathrm{S},\mathrm{R},\mathrm{P}\rangle,
    \qquad \ell=3,
\]
\texttt{RoE} 또는 \texttt{RoW}를 사용하면
\[
    \tau_{\mathrm{NS}}(s)=\langle \mathrm{P},\mathrm{S},\mathrm{R}\rangle,
    \qquad \ell=5,
\]
\texttt{RoL} 또는 \texttt{RoR}을 사용하면
\[
    \tau_{\mathrm{UD}}(s)=\langle \mathrm{R},\mathrm{P},\mathrm{S}\rangle,
    \qquad \ell=4.
\]
같은 말에서도 Rotation 축에 따라 바로 다음 이동거리가 달라진다는 점이 핵심이다.

이하의 도해는 빈 보드에서 말 하나만 놓였다고 가정하며, 시작 Gesture State는 Wrist Axis가 $\mathrm{NS}$인 기준 상태
\[
    \mathrm{S}:\langle \mathrm{S},\mathrm{R},\mathrm{P}\rangle,\qquad
    \mathrm{R}:\langle \mathrm{R},\mathrm{S},\mathrm{P}\rangle,\qquad
    \mathrm{P}:\langle \mathrm{P},\mathrm{R},\mathrm{S}\rangle
\]
를 사용한다. 실제 말의 Wrist Axis가 $\mathrm{EW}$이면 먼저 그 말의 Gesture State를 위 표에서 읽고 같은 연산으로 계산한다. 가운데의 옅은 음영 칸이 시작칸이며, 도달 가능한 칸에는 가능한 최종 윗면을 bar 없이 \texttt{S}, \texttt{R}, \texttt{P}로 적는다. 한 칸에 여러 글자가 적혀 있으면 서로 다른 합법 경로를 통해 여러 윗면이 가능하다는 뜻이고, 순서는 항상 \texttt{S} $\rightarrow$ \texttt{R} $\rightarrow$ \texttt{P}이다.

\section{Basic Movement}

\medskip
\noindent\textbf{Scissors.}
Scissors는 기본 이동거리 3을 가진다.
\input{Figures/Movement/basic-scissors.tex}

\medskip
\noindent\textbf{Rock.}
Rock은 기본 이동거리 4를 가진다. 네 번의 Roll로 작은 고리를 만들 수 있으므로 시작칸으로 되돌아오는 이동도 가능하다.
\input{Figures/Movement/basic-rock.tex}

\medskip
\noindent\textbf{Paper.}
Paper는 기본 이동거리 5를 가지며 세 말 가운데 가장 넓은 범위를 다룬다.
\begin{figure}[htbp]
    \centering
    \begin{tikzpicture}[x=0.620cm,y=0.620cm]
        \fill[black!4] (5,5) rectangle (6,6);
        \draw[rpsc board line] (0,0)--(0,11);
        \draw[rpsc board line] (0,0)--(11,0);
        \draw[rpsc board line] (1,0)--(1,11);
        \draw[rpsc board line] (0,1)--(11,1);
        \draw[rpsc board line] (2,0)--(2,11);
        \draw[rpsc board line] (0,2)--(11,2);
        \draw[rpsc board line] (3,0)--(3,11);
        \draw[rpsc board line] (0,3)--(11,3);
        \draw[rpsc board line] (4,0)--(4,11);
        \draw[rpsc board line] (0,4)--(11,4);
        \draw[rpsc board line] (5,0)--(5,11);
        \draw[rpsc board line] (0,5)--(11,5);
        \draw[rpsc board line] (6,0)--(6,11);
        \draw[rpsc board line] (0,6)--(11,6);
        \draw[rpsc board line] (7,0)--(7,11);
        \draw[rpsc board line] (0,7)--(11,7);
        \draw[rpsc board line] (8,0)--(8,11);
        \draw[rpsc board line] (0,8)--(11,8);
        \draw[rpsc board line] (9,0)--(9,11);
        \draw[rpsc board line] (0,9)--(11,9);
        \draw[rpsc board line] (10,0)--(10,11);
        \draw[rpsc board line] (0,10)--(11,10);
        \draw[rpsc board line] (11,0)--(11,11);
        \draw[rpsc board line] (0,11)--(11,11);
        \draw[rpsc board border] (0,0) rectangle (11,11);
        \node[rpsc analysis text] at (5.50,10.50) {R};
        \node[rpsc analysis text] at (4.50,9.50) {SP};
        \node[rpsc analysis text] at (6.50,9.50) {SP};
        \node[rpsc analysis text compact] at (3.50,8.50) {SRP};
        \node[rpsc analysis text compact] at (5.50,8.50) {SRP};
        \node[rpsc analysis text compact] at (7.50,8.50) {SRP};
        \node[rpsc analysis text compact] at (2.50,7.50) {SRP};
        \node[rpsc analysis text compact] at (4.50,7.50) {SRP};
        \node[rpsc analysis text compact] at (6.50,7.50) {SRP};
        \node[rpsc analysis text compact] at (8.50,7.50) {SRP};
        \node[rpsc analysis text] at (1.50,6.50) {RP};
        \node[rpsc analysis text compact] at (3.50,6.50) {SRP};
        \node[rpsc analysis text compact] at (5.50,6.50) {SRP};
        \node[rpsc analysis text compact] at (7.50,6.50) {SRP};
        \node[rpsc analysis text] at (9.50,6.50) {RP};
        \node[rpsc analysis text] at (0.50,5.50) {S};
        \node[rpsc analysis text compact] at (2.50,5.50) {SRP};
        \node[rpsc analysis text compact] at (4.50,5.50) {SRP};
        \node[rpsc analysis text compact] at (6.50,5.50) {SRP};
        \node[rpsc analysis text compact] at (8.50,5.50) {SRP};
        \node[rpsc analysis text] at (10.50,5.50) {S};
        \node[rpsc analysis text] at (1.50,4.50) {RP};
        \node[rpsc analysis text compact] at (3.50,4.50) {SRP};
        \node[rpsc analysis text compact] at (5.50,4.50) {SRP};
        \node[rpsc analysis text compact] at (7.50,4.50) {SRP};
        \node[rpsc analysis text] at (9.50,4.50) {RP};
        \node[rpsc analysis text compact] at (2.50,3.50) {SRP};
        \node[rpsc analysis text compact] at (4.50,3.50) {SRP};
        \node[rpsc analysis text compact] at (6.50,3.50) {SRP};
        \node[rpsc analysis text compact] at (8.50,3.50) {SRP};
        \node[rpsc analysis text compact] at (3.50,2.50) {SRP};
        \node[rpsc analysis text compact] at (5.50,2.50) {SRP};
        \node[rpsc analysis text compact] at (7.50,2.50) {SRP};
        \node[rpsc analysis text] at (4.50,1.50) {SP};
        \node[rpsc analysis text] at (6.50,1.50) {SP};
        \node[rpsc analysis text] at (5.50,0.50) {R};
    \end{tikzpicture}
    \caption{Paper}
    \label{fig:paper-basic}
\end{figure}


\section{Movement with Items}

Push 도해는 Push 후 첫 Roll부터 즉시 역행 금지를 적용하여 계산한다. Rotation 도해는 여섯 Rotation을 각각 먼저 적용하고, Rotation 후 Top Gesture에 따른 3/4/5 Roll의 모든 합법 경로를 합쳐 표시한다. Step 도해는 Step Short와 Step Long의 합법 경로를 함께 표시한다.

\input{Figures/Movement/scissors-push.tex}
\begin{figure}[htbp]
    \centering
    \begin{tikzpicture}[x=0.620cm,y=0.620cm]
        \fill[black!4] (3,3) rectangle (4,4);
        \foreach \i in {0,...,7} {\draw[rpsc board line] (\i,0)--(\i,7);\draw[rpsc board line] (0,\i)--(7,\i);}
        \draw[rpsc board border] (0,0) rectangle (7,7);
        \node[rpsc analysis text] at (3.50,6.50) {P};
        \node[rpsc analysis text] at (2.50,5.50) {SR};
        \node[rpsc analysis text] at (4.50,5.50) {SR};
        \node[rpsc analysis text] at (1.50,4.50) {SP};
        \node[rpsc analysis text] at (3.50,4.50) {S};
        \node[rpsc analysis text] at (5.50,4.50) {SP};
        \node[rpsc analysis text] at (0.50,3.50) {R};
        \node[rpsc analysis text] at (2.50,3.50) {S};
        \node[rpsc analysis text] at (4.50,3.50) {S};
        \node[rpsc analysis text] at (6.50,3.50) {R};
        \node[rpsc analysis text] at (1.50,2.50) {SP};
        \node[rpsc analysis text] at (3.50,2.50) {S};
        \node[rpsc analysis text] at (5.50,2.50) {SP};
        \node[rpsc analysis text] at (2.50,1.50) {SR};
        \node[rpsc analysis text] at (4.50,1.50) {SR};
        \node[rpsc analysis text] at (3.50,0.50) {P};
    \end{tikzpicture}
    \caption{Scissors with Rotation}
    \label{fig:scissors-rotation}
\end{figure}

\begin{figure}[htbp]
    \centering
    \begin{tikzpicture}[x=0.620cm,y=0.620cm]
        \fill[black!4] (4,4) rectangle (5,5);
        \draw[rpsc board line] (0,0)--(0,9);
        \draw[rpsc board line] (0,0)--(9,0);
        \draw[rpsc board line] (1,0)--(1,9);
        \draw[rpsc board line] (0,1)--(9,1);
        \draw[rpsc board line] (2,0)--(2,9);
        \draw[rpsc board line] (0,2)--(9,2);
        \draw[rpsc board line] (3,0)--(3,9);
        \draw[rpsc board line] (0,3)--(9,3);
        \draw[rpsc board line] (4,0)--(4,9);
        \draw[rpsc board line] (0,4)--(9,4);
        \draw[rpsc board line] (5,0)--(5,9);
        \draw[rpsc board line] (0,5)--(9,5);
        \draw[rpsc board line] (6,0)--(6,9);
        \draw[rpsc board line] (0,6)--(9,6);
        \draw[rpsc board line] (7,0)--(7,9);
        \draw[rpsc board line] (0,7)--(9,7);
        \draw[rpsc board line] (8,0)--(8,9);
        \draw[rpsc board line] (0,8)--(9,8);
        \draw[rpsc board line] (9,0)--(9,9);
        \draw[rpsc board line] (0,9)--(9,9);
        \draw[rpsc board border] (0,0) rectangle (9,9);
        \node[rpsc analysis text] at (4.50,8.50) {S};
        \node[rpsc analysis text] at (3.50,7.50) {RP};
        \node[rpsc analysis text] at (5.50,7.50) {RP};
        \node[rpsc analysis text compact] at (2.50,6.50) {SRP};
        \node[rpsc analysis text compact] at (4.50,6.50) {SRP};
        \node[rpsc analysis text compact] at (6.50,6.50) {SRP};
        \node[rpsc analysis text] at (1.50,5.50) {RP};
        \node[rpsc analysis text] at (3.50,5.50) {RP};
        \node[rpsc analysis text] at (5.50,5.50) {RP};
        \node[rpsc analysis text] at (7.50,5.50) {RP};
        \node[rpsc analysis text] at (0.50,4.50) {S};
        \node[rpsc analysis text compact] at (2.50,4.50) {SRP};
        \node[rpsc analysis text] at (4.50,4.50) {RP};
        \node[rpsc analysis text compact] at (6.50,4.50) {SRP};
        \node[rpsc analysis text] at (8.50,4.50) {S};
        \node[rpsc analysis text] at (1.50,3.50) {RP};
        \node[rpsc analysis text] at (3.50,3.50) {RP};
        \node[rpsc analysis text] at (5.50,3.50) {RP};
        \node[rpsc analysis text] at (7.50,3.50) {RP};
        \node[rpsc analysis text compact] at (2.50,2.50) {SRP};
        \node[rpsc analysis text compact] at (4.50,2.50) {SRP};
        \node[rpsc analysis text compact] at (6.50,2.50) {SRP};
        \node[rpsc analysis text] at (3.50,1.50) {RP};
        \node[rpsc analysis text] at (5.50,1.50) {RP};
        \node[rpsc analysis text] at (4.50,0.50) {S};
    \end{tikzpicture}
    \caption{Scissors with Step}
    \label{fig:scissors-step}
\end{figure}


\begin{figure}[htbp]
    \centering
    \begin{tikzpicture}[x=0.620cm,y=0.620cm]
        \fill[black!4] (5,5) rectangle (6,6);
        \foreach \i in {0,...,11} {\draw[rpsc board line] (\i,0)--(\i,11);\draw[rpsc board line] (0,\i)--(11,\i);}
        \draw[rpsc board border] (0,0) rectangle (11,11);
        \node[rpsc analysis text] at (5.50,10.50) {R};
        \node[rpsc analysis text compact] at (4.50,9.50) {SRP};\node[rpsc analysis text compact] at (6.50,9.50) {SRP};
        \node[rpsc analysis text compact] at (3.50,8.50) {SRP};\node[rpsc analysis text compact] at (5.50,8.50) {SRP};\node[rpsc analysis text compact] at (7.50,8.50) {SRP};
        \node[rpsc analysis text compact] at (2.50,7.50) {SRP};\node[rpsc analysis text compact] at (4.50,7.50) {SRP};\node[rpsc analysis text compact] at (6.50,7.50) {SRP};\node[rpsc analysis text compact] at (8.50,7.50) {SRP};
        \node[rpsc analysis text compact] at (1.50,6.50) {SRP};\node[rpsc analysis text compact] at (3.50,6.50) {SRP};\node[rpsc analysis text compact] at (5.50,6.50) {SRP};\node[rpsc analysis text compact] at (7.50,6.50) {SRP};\node[rpsc analysis text compact] at (9.50,6.50) {SRP};
        \node[rpsc analysis text] at (0.50,5.50) {R};\node[rpsc analysis text compact] at (2.50,5.50) {SRP};\node[rpsc analysis text compact] at (4.50,5.50) {SRP};\node[rpsc analysis text compact] at (6.50,5.50) {SRP};\node[rpsc analysis text compact] at (8.50,5.50) {SRP};\node[rpsc analysis text] at (10.50,5.50) {R};
        \node[rpsc analysis text compact] at (1.50,4.50) {SRP};\node[rpsc analysis text compact] at (3.50,4.50) {SRP};\node[rpsc analysis text compact] at (5.50,4.50) {SRP};\node[rpsc analysis text compact] at (7.50,4.50) {SRP};\node[rpsc analysis text compact] at (9.50,4.50) {SRP};
        \node[rpsc analysis text compact] at (2.50,3.50) {SRP};\node[rpsc analysis text compact] at (4.50,3.50) {SRP};\node[rpsc analysis text compact] at (6.50,3.50) {SRP};\node[rpsc analysis text compact] at (8.50,3.50) {SRP};
        \node[rpsc analysis text compact] at (3.50,2.50) {SRP};\node[rpsc analysis text compact] at (5.50,2.50) {SRP};\node[rpsc analysis text compact] at (7.50,2.50) {SRP};
        \node[rpsc analysis text compact] at (4.50,1.50) {SRP};\node[rpsc analysis text compact] at (6.50,1.50) {SRP};
        \node[rpsc analysis text] at (5.50,0.50) {R};
    \end{tikzpicture}
    \caption{Rock with Push}
    \label{fig:rock-push}
\end{figure}

\begin{figure}[htbp]
    \centering
    \begin{tikzpicture}[x=0.620cm,y=0.620cm]
        \fill[black!4] (4,4) rectangle (5,5);
        \foreach \i in {0,...,9} {\draw[rpsc board line] (\i,0)--(\i,9);\draw[rpsc board line] (0,\i)--(9,\i);}
        \draw[rpsc board border] (0,0) rectangle (9,9);
        \node[rpsc analysis text] at (4.50,8.50) {R};
        \node[rpsc analysis text] at (3.50,7.50) {SP};\node[rpsc analysis text] at (5.50,7.50) {SP};
        \node[rpsc analysis text compact] at (2.50,6.50) {SRP};\node[rpsc analysis text compact] at (4.50,6.50) {SRP};\node[rpsc analysis text compact] at (6.50,6.50) {SRP};
        \node[rpsc analysis text] at (1.50,5.50) {SP};\node[rpsc analysis text] at (3.50,5.50) {SP};\node[rpsc analysis text] at (5.50,5.50) {SP};\node[rpsc analysis text] at (7.50,5.50) {SP};
        \node[rpsc analysis text] at (0.50,4.50) {R};\node[rpsc analysis text compact] at (2.50,4.50) {SRP};\node[rpsc analysis text] at (4.50,4.50) {SP};\node[rpsc analysis text compact] at (6.50,4.50) {SRP};\node[rpsc analysis text] at (8.50,4.50) {R};
        \node[rpsc analysis text] at (1.50,3.50) {SP};\node[rpsc analysis text] at (3.50,3.50) {SP};\node[rpsc analysis text] at (5.50,3.50) {SP};\node[rpsc analysis text] at (7.50,3.50) {SP};
        \node[rpsc analysis text compact] at (2.50,2.50) {SRP};\node[rpsc analysis text compact] at (4.50,2.50) {SRP};\node[rpsc analysis text compact] at (6.50,2.50) {SRP};
        \node[rpsc analysis text] at (3.50,1.50) {SP};\node[rpsc analysis text] at (5.50,1.50) {SP};
        \node[rpsc analysis text] at (4.50,0.50) {R};
    \end{tikzpicture}
    \caption{Rock with Rotation}
    \label{fig:rock-rotation}
\end{figure}

\input{Figures/Movement/rock-step.tex}

\begin{figure}[htbp]
    \centering
    \begin{tikzpicture}[x=0.620cm,y=0.620cm]
        \fill[black!4] (6,6) rectangle (7,7);
        \draw[rpsc board line] (0,0)--(0,13);
        \draw[rpsc board line] (0,0)--(13,0);
        \draw[rpsc board line] (1,0)--(1,13);
        \draw[rpsc board line] (0,1)--(13,1);
        \draw[rpsc board line] (2,0)--(2,13);
        \draw[rpsc board line] (0,2)--(13,2);
        \draw[rpsc board line] (3,0)--(3,13);
        \draw[rpsc board line] (0,3)--(13,3);
        \draw[rpsc board line] (4,0)--(4,13);
        \draw[rpsc board line] (0,4)--(13,4);
        \draw[rpsc board line] (5,0)--(5,13);
        \draw[rpsc board line] (0,5)--(13,5);
        \draw[rpsc board line] (6,0)--(6,13);
        \draw[rpsc board line] (0,6)--(13,6);
        \draw[rpsc board line] (7,0)--(7,13);
        \draw[rpsc board line] (0,7)--(13,7);
        \draw[rpsc board line] (8,0)--(8,13);
        \draw[rpsc board line] (0,8)--(13,8);
        \draw[rpsc board line] (9,0)--(9,13);
        \draw[rpsc board line] (0,9)--(13,9);
        \draw[rpsc board line] (10,0)--(10,13);
        \draw[rpsc board line] (0,10)--(13,10);
        \draw[rpsc board line] (11,0)--(11,13);
        \draw[rpsc board line] (0,11)--(13,11);
        \draw[rpsc board line] (12,0)--(12,13);
        \draw[rpsc board line] (0,12)--(13,12);
        \draw[rpsc board line] (13,0)--(13,13);
        \draw[rpsc board line] (0,13)--(13,13);
        \draw[rpsc board border] (0,0) rectangle (13,13);
        \node[rpsc analysis text] at (6.50,12.50) {R};
        \node[rpsc analysis text compact] at (5.50,11.50) {SRP};
        \node[rpsc analysis text compact] at (7.50,11.50) {SRP};
        \node[rpsc analysis text compact] at (4.50,10.50) {SRP};
        \node[rpsc analysis text compact] at (6.50,10.50) {SRP};
        \node[rpsc analysis text compact] at (8.50,10.50) {SRP};
        \node[rpsc analysis text compact] at (3.50,9.50) {SRP};
        \node[rpsc analysis text compact] at (5.50,9.50) {SRP};
        \node[rpsc analysis text compact] at (7.50,9.50) {SRP};
        \node[rpsc analysis text compact] at (9.50,9.50) {SRP};
        \node[rpsc analysis text compact] at (2.50,8.50) {SRP};
        \node[rpsc analysis text compact] at (4.50,8.50) {SRP};
        \node[rpsc analysis text compact] at (6.50,8.50) {SRP};
        \node[rpsc analysis text compact] at (8.50,8.50) {SRP};
        \node[rpsc analysis text compact] at (10.50,8.50) {SRP};
        \node[rpsc analysis text compact] at (1.50,7.50) {SRP};
        \node[rpsc analysis text compact] at (3.50,7.50) {SRP};
        \node[rpsc analysis text compact] at (5.50,7.50) {SRP};
        \node[rpsc analysis text compact] at (7.50,7.50) {SRP};
        \node[rpsc analysis text compact] at (9.50,7.50) {SRP};
        \node[rpsc analysis text compact] at (11.50,7.50) {SRP};
        \node[rpsc analysis text] at (0.50,6.50) {S};
        \node[rpsc analysis text compact] at (2.50,6.50) {SRP};
        \node[rpsc analysis text compact] at (4.50,6.50) {SRP};
        \node[rpsc analysis text compact] at (6.50,6.50) {SRP};
        \node[rpsc analysis text compact] at (8.50,6.50) {SRP};
        \node[rpsc analysis text compact] at (10.50,6.50) {SRP};
        \node[rpsc analysis text] at (12.50,6.50) {S};
        \node[rpsc analysis text compact] at (1.50,5.50) {SRP};
        \node[rpsc analysis text compact] at (3.50,5.50) {SRP};
        \node[rpsc analysis text compact] at (5.50,5.50) {SRP};
        \node[rpsc analysis text compact] at (7.50,5.50) {SRP};
        \node[rpsc analysis text compact] at (9.50,5.50) {SRP};
        \node[rpsc analysis text compact] at (11.50,5.50) {SRP};
        \node[rpsc analysis text compact] at (2.50,4.50) {SRP};
        \node[rpsc analysis text compact] at (4.50,4.50) {SRP};
        \node[rpsc analysis text compact] at (6.50,4.50) {SRP};
        \node[rpsc analysis text compact] at (8.50,4.50) {SRP};
        \node[rpsc analysis text compact] at (10.50,4.50) {SRP};
        \node[rpsc analysis text compact] at (3.50,3.50) {SRP};
        \node[rpsc analysis text compact] at (5.50,3.50) {SRP};
        \node[rpsc analysis text compact] at (7.50,3.50) {SRP};
        \node[rpsc analysis text compact] at (9.50,3.50) {SRP};
        \node[rpsc analysis text compact] at (4.50,2.50) {SRP};
        \node[rpsc analysis text compact] at (6.50,2.50) {SRP};
        \node[rpsc analysis text compact] at (8.50,2.50) {SRP};
        \node[rpsc analysis text compact] at (5.50,1.50) {SRP};
        \node[rpsc analysis text compact] at (7.50,1.50) {SRP};
        \node[rpsc analysis text] at (6.50,0.50) {R};
    \end{tikzpicture}
    \caption{Paper with Push}
    \label{fig:paper-push}
\end{figure}

\input{Figures/Movement/paper-rotation.tex}
\input{Figures/Movement/paper-step.tex}
